\scanpage{104}
Variante globale : on définit des hypertor globaux \unsure{en} prenant
\[
\mathbb R\Gamma_X\ L_\bullet\overset{\mathbb L}{\otimes}M_\bullet.
\]
\margin{Mais il faut supposer, \unsure{si on ne fait pas d'hyp.} sur $L,M$, que $\Gamma_X$ est de dim \unsure{coh} finie !}

\uline{Lois formelles les plus évidentes} (lois internes).

1) Associativité de $\overset{\mathbb L}{\otimes}$

2) Commutativité (attention au signe !)

3) \diagram{\begin{tikzcd}[column sep=small,row sep=small]
\mathbb R\mathcal Hom(L_\bullet\overset{\mathbb L}{\otimes}M_\bullet,K^\bullet) & \simeq\mathbb R\mathcal Hom(L_\bullet,\mathbb R\mathcal Hom(M_\bullet,K^\bullet)) \\
& \simeq\mathbb R\mathcal Hom(M_\bullet,\mathbb R\mathcal Hom(L_\bullet,K^\bullet)) \\
\mathbb R\operatorname{Hom}(L_\bullet\overset{\mathbb L}{\otimes}M_\bullet,K^\bullet)\arrow[uu,bend left=65] & \simeq\mathbb R\operatorname{Hom}(L_\bullet,\mathbb R\mathcal Hom(M_\bullet,K^\bullet)) \\
& \simeq\mathbb R\operatorname{Hom}(M_\bullet,\mathbb R\mathcal Hom(L_\bullet,K^\bullet))
\end{tikzcd}}
\margin{N.B. $\mathbb R\mathcal Hom(L_\bullet,K^\bullet)$ défini dès que $K^\bullet\in\operatorname{Ob}D^+(X)$, mais $L_\bullet$ peut être quelconque}
\[
\deleted{\operatorname{Hom}(L_\bullet,\mathbb R\mathcal Hom}
\]
\[
\begin{aligned}
\operatorname{Hom}(L_\bullet\overset{\mathbb L}{\otimes}M_\bullet,K^\bullet)&\simeq\operatorname{Hom}(L_\bullet,\mathbb R\mathcal Hom(M_\bullet,K^\bullet)\\
&\simeq\operatorname{Hom}(M_\bullet,\mathbb R\mathcal Hom(L_\bullet,K^\bullet))
\end{aligned}
\]

\uline{Remarque sur changement de base non plat}

On définit $\mathbb Lf^*:\unsure{D^-|D_0^-}(Y)\to D^-(X)$ encore en passant par $K^-(Y)^{\mathrm{plat}}/\text{quasi-isom}$, en notant que $f^*$ \deleted{sur} \deleted{$\unsure{K^-}$} \unsure{mène} $K^-(Y)^{\mathrm{plat}}$ dans $K^-(X)^{\mathrm{plat}}$ et transforme quasi. isom en \unsure{itou.|idem.} Avec cette définition, on aura encore
\[
\mathbb Lf^*(L_\bullet\overset{\mathbb L}{\otimes}M_\bullet)\simeq\mathbb Lf^*(L_\bullet)\overset{\mathbb L}{\otimes}\mathbb Lf^*(M_\bullet)
\]
Notons aussi formule de transitivité
\mspage{103}
